\begin{exercise}{89}
Show that a Joyal equivalence $f\colon\mathcal C\to\mathcal D$ induces an
equivalence of ordinary categories $h\mathcal C\to h\mathcal D$.

\begin{proof}
By Remark~2.1.15, choose $g\colon\mathcal D\to\mathcal C$ with natural
equivalences $gf\simeq\mathrm{id}_{\mathcal C}$ and
$fg\simeq\mathrm{id}_{\mathcal D}$.
For any natural equivalence $\tau\colon u\to v$ between functors
$\mathcal A\to\mathcal B$, its adjoint
$H\colon\mathcal A\times\Delta^1\to\mathcal B$ induces a functor
$h\mathcal A\times[1]\to h\mathcal B$ by Exercise~\ref{ex:28}
and Corollary~1.2.14.
This is a natural transformation $hu\to hv$.
Its components are isomorphisms: evaluation at each object sends
the equivalence $\tau$ to an equivalence of $\mathcal B$
(Remark~2.1.17 and Definition~1.2.22).
Applying this to the chosen natural equivalences gives natural isomorphisms
$(hg)(hf)\cong\mathrm{id}_{h\mathcal C}$ and
$(hf)(hg)\cong\mathrm{id}_{h\mathcal D}$.
Thus $hf$ is an equivalence with quasi-inverse $hg$.
\end{proof}
\end{exercise}
